% ===== 5. 議論と限界 (Discussion & Limitations) =====
% 役割:
%   - Self-Audit Principle の位置付け (locked-in: ここに置く、名称は弱める)
%   - 本研究の限界を率直に列挙 → 査読者の懸念を先回りで回収
%
% 分量目安: 1 ページ
%
% locked-in 反映:
%   - Self-Audit Principle → 「出力の意味的誠実性に対する自己整合性評価」
%     という neutral な表現で言及
%   - 「無意識的重力仮説 (UGH)」は理論名として 1 段落のみ言及

\section{議論}
\label{sec:discussion}

% TODO: 5.1 自己整合性評価への適用可能性
% Self-Audit Principle を「出力の意味的誠実性に対する自己整合性評価」
% として再フレーミング。LLM 出力の自己診断ループとしての応用を議論。
\subsection{自己整合性評価への適用可能性}
% [PLACEHOLDER]

% TODO: 5.2 理論的位置付け
% 本指標群は著者らの提唱する「無意識的重力仮説 (UGH)」の枠組みに
% 位置付けられる。詳細な理論的動機は別稿に譲る。
\subsection{理論的位置付け}
% [PLACEHOLDER - 1 段落のみ]

% ----- 限界 -----
\section{限界}
\label{sec:limitations}

% TODO: 5.3 限界 (率直に列挙)
\subsection{本研究の限界}
% [PLACEHOLDER]
% - サンプル規模: HA48 (拡充後 HA100) は予備検証規模
% - annotator 数: 単一 annotator (拡充計画で第二 annotator + IAA を追加予定)
% - 単一ドメイン: AI 応答監査タスクに限定、汎化性は今後の検証課題
% - ベースライン: G-Eval / BERTScore のみ、より広範な比較は今後
% - 閾値固定性: HA48 で校正、拡充版での再校正は行わない (再現性確認のため)
